% =============================================================================
% Fichier  : chord-substitution-fr.tex
% Titre    : La substitution par corde clôt l'équation de Jensen à coefficient continu
% Sous-titre : Une note pédagogique sur trois hypothèses de régularité vestigiales
% Auteur   : [AUTEUR]
% Revue    : Comptes Rendus Mathématique (Académie des sciences)
% Source   : satellites/o3-chord-substitution/13-draft0.4-fr-manuscrit.md
% Date     : 2026-06-06
%
% Compilation :
%   pdflatex chord-substitution-fr
%   bibtex   chord-substitution-fr
%   pdflatex chord-substitution-fr
%   pdflatex chord-substitution-fr
%
% Classe   : La classe Centre Mersenne `crmath` est le modèle officiel
%            (télécharger pack_author-crmath.zip depuis
%            centre-mersenne.org/pack_author/).
%            Une configuration de repli amsart est fournie ci-dessous pour
%            les utilisateurs qui n'ont pas la classe crmath installée :
%            commenter \documentclass{crmath} et décommenter le bloc amsart.
% =============================================================================

% --- CLASSE CRAS OFFICIELLE (préférée) ---
% \documentclass{crmath}

% --- CONFIGURATION DE REPLI AMSART (active par défaut) ---
\documentclass[11pt,a4paper]{amsart}
\usepackage[utf8]{inputenc}
\usepackage[T1]{fontenc}
\usepackage{amssymb,amsmath,amsthm}
\usepackage{microtype}
\usepackage{booktabs}
\usepackage{array}
\usepackage{tabularx}
\usepackage[hidelinks]{hyperref}
\usepackage[english,french]{babel} % `french` actif par défaut, espace insécable automatique avant :;?!
\usepackage{csquotes}
\usepackage{lineno}
\linenumbers
\usepackage[a4paper,margin=2.5cm]{geometry}
\renewcommand{\baselinestretch}{1.5} % double interligne (exigence CRAS)
\setlength{\parindent}{1em}
\setlength{\parskip}{0.25em}

% Numérotation des sections en chiffres romains (convention CRAS Mathématique)
\renewcommand{\thesection}{\Roman{section}}
\renewcommand{\thesubsection}{\Roman{section}.\arabic{subsection}}

% Environnements de théorème (en français)
\newtheorem{theorem}{Théorème}
\newtheorem{corollary}[theorem]{Corollaire}
\newtheorem{proposition}[theorem]{Proposition}
\newtheorem{lemma}[theorem]{Lemme}
\newtheorem{remark}[theorem]{Remarque}

% Renommage français des étiquettes auto-générées
\renewcommand{\proofname}{Démonstration}
\renewcommand{\abstractname}{Résumé}
\renewcommand{\refname}{Références bibliographiques}

% Commandes personnalisées
\newcommand{\R}{\mathbb R}
\newcommand{\Q}{\mathbb Q}
\newcommand{\N}{\mathbb N}
\newcommand{\E}{\mathbb E}
\newcommand{\PR}{\mathbb P}
% --- FIN DU REPLI ---

\begin{document}

\title[La substitution par corde clôt l'équation de Jensen à coefficient continu]%
{La substitution par corde clôt l'équation de Jensen à coefficient continu\\
\normalsize\textit{Une note pédagogique sur trois hypothèses de régularité vestigiales}}

\author{[AUTEUR]\thanks{Auteur correspondant. Affiliation :
[INSTITUTION, PAYS]. Courriel : [auteur@institution.pays].}}

\date{\today}

\subjclass[2020]{39B22, 39B05}
\keywords{équation de Jensen, équation de Cauchy, base de Hamel, équation
fonctionnelle, fonction affine, substitution par corde}

\begin{abstract}
La forme à coefficient continu de l'équation fonctionnelle de Jensen sur
un intervalle réel --- celle où le coefficient parcourt tout le
continuum~$[0,1]$ plutôt qu'une valeur unique telle que $\tfrac12$ ---
ne partage aucune pathologie de type base de Hamel avec son analogue à
coefficient discret, et se résout en forme close par une substitution par
corde tenant en une ligne.

Plus précisément~: pour $G\colon[0,M]\to\R$ satisfaisant
$p\,G(u_1) + (1-p)\,G(u_2) = G(p\,u_1 + (1-p)\,u_2)$ pour tous
$u_1,u_2\in[0,M]$ et tout $p\in[0,1]$, la fonction $G$ est affine
sur~$[0,M]$ --- c.-à-d. de la forme $G(v) = av + b$ avec $a,b\in\R$
--- sans qu'aucune hypothèse de mesurabilité, de bornitude ou de
continuité ne soit requise.

Nous en donnons la démonstration, exhibons la formule affine explicite,
articulons le dictionnaire des trois hypothèses de régularité
classiquement requises qui deviennent vestigiales sous la forme à
coefficient continu, et documentons une occurrence du piège dans la
littérature sur la calibration par perte de substitution, ainsi que la
raison structurelle pour laquelle il est appelé à se répéter dans toute
dérivation qui pousse l'inégalité de Jensen à saturation sur une classe
assez large de distributions bipoints.
\end{abstract}

% Métadonnées bilingues : version anglaise du titre, du résumé et des
% mots-clés. Le manuscrit anglais détaillé est dans chord-substitution-en.tex
% (fichier séparé).
\renewcommand{\abstractname}{Abstract}
\begin{abstract}
\selectlanguage{english}
The continuous-coefficient form of the Jensen functional equation on a
real interval --- the form in which the coefficient varies over the full
continuum $[0,1]$ rather than over a single value such as $\tfrac12$
--- shares no Hamel-basis pathology with its discrete-coefficient sibling,
and is solved in closed form by a one-line chord substitution. Concretely:
for $G\colon[0,M]\to\R$ satisfying $pG(u_1)+(1-p)G(u_2)=G(pu_1+(1-p)u_2)$
for all $u_1,u_2\in[0,M]$ and all $p\in[0,1]$, $G$ is affine on $[0,M]$
--- meaning $G(v)=av+b$ for some $a,b\in\R$ --- with no measurability,
boundedness, or continuity hypothesis on $G$. We give the proof, exhibit
the explicit affine formula, articulate the dictionary of three classically
required regularity hypotheses that become vestigial under the
continuous-coefficient form, and document one recurrence of the trap in the
surrogate-calibration literature.

\textbf{Keywords:} Jensen equation; Cauchy equation; Hamel basis;
functional equation; affine function; chord substitution.
\selectlanguage{french}
\end{abstract}
\renewcommand{\abstractname}{Résumé}

\maketitle

% =============================================================================
\section{Introduction}
\label{sec:introduction}

L'équation
\begin{equation}
p\,G(u_1) + (1-p)\,G(u_2) \;=\; G\bigl(p\,u_1 + (1-p)\,u_2\bigr)
\qquad (u_1,u_2\in I,\ p\in[0,1])
\tag{$\star$}\label{eq:star}
\end{equation}
sur un intervalle réel~$I$ --- où $G\colon I\to\R$ est l'inconnue et
$u_1,u_2,p$ sont les arguments et le coefficient --- a au moins trois
visages en analyse classique, et ces visages ont été confondus sans relâche
depuis plus d'un siècle.

L'\textbf{équation de Jensen à coefficient discret},
\begin{equation}
G\!\left(\tfrac{u_1 + u_2}{2}\right) \;=\;
\tfrac{G(u_1) + G(u_2)}{2},
\tag{$J_2$}\label{eq:J2}
\end{equation}
est la spécialisation de~\eqref{eq:star} au cas $p = \tfrac12$. En posant
$f(x) := G(x) - G(0)$, l'équation~\eqref{eq:J2} équivaut, sur un translaté
de~$I$, à l'équation additive de Cauchy $f(x + y) = f(x) + f(y)$
(cf.~\cite{Aczel1966} ou~\cite{Kuczma2009}). L'additivité de Cauchy hérite
alors d'un appareil complet de solutions pathologiques~: en l'absence de
mesurabilité, de monotonie, de bornitude sur un ensemble de mesure positive
ou de toute autre hypothèse de régularité, l'équation admet des solutions
non affines construites au moyen d'une \textbf{base de Hamel} --- c.-à-d.
d'une base de~$\R$ en tant que $\Q$-espace vectoriel, dont l'existence
requiert l'axiome du choix --- voir~\cite{Hamel1905}. Les théorèmes
classiques de Cauchy~\cite{Cauchy1821}, Darboux~\cite{Darboux1875},
Hamel~\cite{Hamel1905}, Ostrowski~\cite{Ostrowski1929},
Sierpiński~\cite{Sierpinski1920} et Steinhaus~\cite{Steinhaus1920}
délimitent ensemble les hypothèses de régularité suffisantes pour récupérer
l'affinité --- \emph{une} hypothèse au moins est véritablement nécessaire,
faute de quoi l'affinité échoue.

L'\textbf{équation de Jensen à coefficient rationnel},
\begin{equation}
p\,G(u_1) + (1-p)\,G(u_2) \;=\; G(p\,u_1 + (1-p)\,u_2)
\qquad (u_1,u_2\in I,\ p\in[0,1]\cap\Q),
\tag{$J_{\Q}$}\label{eq:JQ}
\end{equation}
raffine~\eqref{eq:J2} sans pour autant échapper à sa pathologie~:
l'additivité de Cauchy entraîne la $\Q$-homogénéité $f(qx) = q\,f(x)$ par
les arguments standard ($G(nx) = nG(x)$ par récurrence, $G(x/n) = G(x)/n$
par substitution, $G(qx) = qG(x)$ pour $q\in\Q$ par combinaison), et la
$\Q$-homogénéité restitue~\eqref{eq:JQ} pour~$f$, donc pour~$G$. Ainsi
\eqref{eq:JQ}, \eqref{eq:J2} et l'équation de Cauchy partagent la même
classe de solutions à une constante près, et toutes trois héritent de la
pathologie de Hamel.

La \textbf{forme à coefficient continu}~\eqref{eq:star}, en revanche, est
fondamentalement différente. C'est l'objet de cette note.

\paragraph{Résultat.}
Supposons que $G\colon[0,M]\to\R$ vérifie~\eqref{eq:star} pour tous
$u_1,u_2\in[0,M]$ et tout $p\in[0,1]$. Alors $G$ est affine sur~$[0,M]$
--- c.-à-d. de la forme $G(v) = av + b$ avec $a,b\in\R$ --- et plus
précisément
\[
G(v) \;=\; a\,v + b,
\qquad a = \frac{G(M) - G(0)}{M},\quad b = G(0).
\]
La démonstration consiste en une unique substitution~: poser $u_1 = M$,
$u_2 = 0$, $p = v/M$ dans~\eqref{eq:star}. Aucune hypothèse de régularité
sur~$G$ n'est consommée.

Il s'agit là d'un résultat folklorique --- \cite{Aczel1966}
et~\cite{Kuczma2009} articulent au moins l'observation sous-jacente, et le
résultat est standard dans la communauté des équations fonctionnelles. Une
confusion durable persiste cependant dans les travaux appliqués qui
rencontrent~\eqref{eq:star} en dehors de la littérature dédiée~: les
auteurs invoquent par réflexe la machinerie de régularité de l'équation de
Cauchy --- continuité, bornitude, mesurabilité --- alors qu'aucune de ces
hypothèses n'intervient effectivement dans la démonstration de l'affinité.
La présente note est une note pédagogique et de citation~: nous exhibons
la démonstration ainsi que le dictionnaire des hypothèses vestigiales
pour~\eqref{eq:star}, et fournissons un point de citation court permettant
aux auteurs qui rencontrent l'équation dans leurs propres travaux de
retirer la préoccupation Hamel sans avoir à la redériver.

La suite est organisée comme suit. La Section~\ref{sec:result} énonce le
théorème et sa démonstration. La Section~\ref{sec:dictionary} est le
dictionnaire des trois hypothèses de régularité qui \emph{ne sont pas}
requises, chacune accompagnée de la pathologie correspondante
pour~\eqref{eq:J2} qui échoue à se manifester pour~\eqref{eq:star}. La
Section~\ref{sec:variants} collecte les variantes --- l'hypothèse minimale
que la démonstration consomme effectivement~; la remarque dimensionnelle~;
le contraste avec~\eqref{eq:JQ}. La Section~\ref{sec:recurrence} documente
une occurrence du piège dans la littérature sur la calibration par perte
de substitution, et énonce la raison structurelle qui rend le piège appelé
à se répéter dans toute dérivation qui pousse l'inégalité de Jensen à
saturation.

% =============================================================================
\section{Le résultat}
\label{sec:result}

\begin{theorem}
\label{thm:main}
Soient $M > 0$ et $G\colon[0,M]\to\R$ vérifiant~\eqref{eq:star} pour tous
$u_1,u_2\in[0,M]$ et tout $p\in[0,1]$. Alors $G$ est affine sur~$[0,M]$~:
\[
G(v) \;=\; \frac{G(M) - G(0)}{M}\,v \;+\; G(0)
\qquad \text{pour tout } v \in [0, M].
\]
Aucune hypothèse de mesurabilité, de bornitude ou de continuité sur~$G$
n'est requise.
\end{theorem}

\begin{proof}
Fixons $v \in [0, M]$. En posant $u_1 = M$, $u_2 = 0$ et $p = v/M \in [0, 1]$
dans~\eqref{eq:star},
\[
\frac{v}{M}\,G(M) + \left(1 - \frac{v}{M}\right) G(0) \;=\;
G\!\left(\frac{v}{M} \cdot M + \left(1 - \frac{v}{M}\right) \cdot 0\right)
\;=\; G(v).
\]
En réarrangeant, $G(v) = G(0) + \dfrac{G(M) - G(0)}{M}\,v$.
\end{proof}

\begin{corollary}
\label{cor:regularity}
Sous les hypothèses du Théorème~\ref{thm:main}, $G$ est en particulier
continue, monotone, localement lipschitzienne, absolument continue et
mesurable sur~$[0, M]$.
\end{corollary}

\begin{proof}
Les fonctions affines vérifient toutes ces propriétés.
\end{proof}

L'ordre d'inférence importe et mérite d'être souligné~: la \emph{conclusion}
du Théorème~\ref{thm:main} est que $G$ possède toutes les propriétés de
régularité que l'on aurait pu vouloir \emph{supposer}. Chacune de ces
propriétés est donc vestigiale en tant qu'hypothèse --- voir la
Section~\ref{sec:dictionary}.

% =============================================================================
\section{Le dictionnaire}
\label{sec:dictionary}

Plaçons côte à côte l'équation à coefficient discret~\eqref{eq:J2} et
l'équation à coefficient continu~\eqref{eq:star}. La première colonne du
Tableau~\ref{tab:dictionary} nomme une hypothèse de régularité
classiquement invoquée pour~\eqref{eq:J2} et en donne une définition en
une phrase~; la deuxième colonne indique si l'hypothèse est requise pour
forcer l'affinité dans le cas~\eqref{eq:J2}, avec l'attribution standard~;
la troisième colonne indique ce qu'il advient de l'hypothèse
sous~\eqref{eq:star}.

\begin{table}[ht]
\renewcommand{\arraystretch}{1.4}
\begin{tabularx}{\textwidth}{@{} X | X | X @{}}
\toprule
\textbf{Hypothèse sur $G$ (définition)} & \textbf{Requise pour
\eqref{eq:J2}$\Rightarrow$ affine~?} & \textbf{Requise pour
\eqref{eq:star}$\Rightarrow$ affine~?}\\
\midrule
\textbf{Continuité sur~$I$} ($G$ est continue en tout $v\in I$).
& Oui, suffisante (Cauchy~\cite{Cauchy1821}).
& \textbf{Non} --- l'affinité est la \emph{conclusion}, non l'hypothèse.\\
\midrule
\textbf{Mesurabilité sur~$I$} ($G$ est mesurable au sens de Lebesgue ou
de Borel).
& Oui, suffisante (Sierpiński~\cite{Sierpinski1920}).
& \textbf{Non} --- idem.\\
\midrule
\textbf{Monotonie sur~$I$} ($G$ est croissante ou décroissante au sens
large).
& Oui, suffisante (Darboux~\cite{Darboux1875}).
& \textbf{Non} --- idem.\\
\midrule
\textbf{Bornitude sur un ensemble de mesure positive} (il existe un
borélien $E\subseteq I$ avec $|E|>0$ sur lequel $G$ est bornée).
& Oui, suffisante (Steinhaus~\cite{Steinhaus1920}~;
cf.~Sierpiński~\cite{Sierpinski1920}).
& \textbf{Non} --- idem.\\
\midrule
\textbf{Bornitude sur~$I$} ($G$ est bornée sur l'intervalle~$I$ entier).
& Oui, suffisante (cas particulier de Steinhaus~\cite{Steinhaus1920}~;
Ostrowski~\cite{Ostrowski1929}).
& \textbf{Non} --- idem.\\
\midrule
\textbf{Aucune} (pas d'hypothèse de régularité du tout).
& \textbf{Insuffisante} --- pathologie de la base de Hamel
(Hamel~\cite{Hamel1905}).
& \textbf{Suffisante} --- Théorème~\ref{thm:main}.\\
\bottomrule
\end{tabularx}
\caption{Hypothèses de régularité sur~$G$. La colonne~2 indique quelles
hypothèses sont classiquement requises pour forcer l'affinité dans le cas
de l'équation à coefficient discret~\eqref{eq:J2}. La colonne~3 indique ce
qu'il en advient sous l'équation à coefficient continu~\eqref{eq:star}.}
\label{tab:dictionary}
\end{table}

Le point essentiel du tableau est sa dernière ligne. \textbf{Aucune
hypothèse de régularité} ne suffit pour l'équation à coefficient discret~:
une fonction additive sur~$\R$, non affine, non mesurable et non bornée
--- la pathologie classique de la base de Hamel de~\cite{Hamel1905} ---
fournit une solution non affine de~\eqref{eq:J2}. Inversement,
\textbf{aucune hypothèse de régularité} n'est \emph{requise} pour
l'équation à coefficient continu~: la substitution par corde du
Théorème~\ref{thm:main} clôt la démonstration au niveau de l'équation
fonctionnelle nue.

Le mécanisme est direct. Une solution pathologique de type base de
Hamel~$G$ de l'équation de Cauchy sur~$\R$ est $\Q$-linéaire~: $G(0) = 0$
et $G(qx) = q\,G(x)$ pour tout $q \in \Q$ et tout $x \in \R$. Par
construction, elle n'est \emph{pas} $\R$-linéaire --- $G(rx) \neq r\,G(x)$
pour un certain irrationnel~$r$ et un certain $x \in \R$.

L'identité de corde à la configuration $u_1 = M$, $u_2 = 0$ s'écrit
\begin{equation}
p\,G(M) + (1 - p)\,G(0) \;=\; G(p\,M).
\tag{$\star_0$}\label{eq:star0}
\end{equation}
Pour~$G$ une fonction additive Hamel-pathologique ($G(0) = 0$),
\eqref{eq:star0} équivaut à l'assertion $p\,G(M) = G(p\,M)$ --- vraie en
$p \in \Q$ (par $\Q$-linéarité) mais fausse aux irrationnels~$p$ où la
$\R$-linéarité est rompue. Ainsi $G$ vérifie~\eqref{eq:JQ} sur $[0, M]$
mais viole~\eqref{eq:star} en \emph{tout} irrationnel $p \in (0, 1)$ où le
défaut de $\R$-linéarité apparaît. L'équation à coefficient
continu~\eqref{eq:star} écarte la pathologie précisément aux points
$p \in [0, 1] \setminus \Q$ sur lesquels les versions à coefficient discret
et rationnel restent muettes. \textbf{La pathologie de Hamel vit aux
irrationnels~$p$ --- précisément là où le coefficient continu
de~\eqref{eq:star} referme la porte.}

% =============================================================================
\section{Variantes et limites}
\label{sec:variants}

\subsection{L'hypothèse minimale que la démonstration consomme}
\label{ssec:weak}

La démonstration du Théorème~\ref{thm:main} consomme~\eqref{eq:star} en une
unique configuration~: $u_1 = M$, $u_2 = 0$, avec $p \in [0, 1]$ libre. La
pleine force de~\eqref{eq:star} --- pour \emph{tout} couple $(u_1, u_2)$
--- n'est pas mobilisée. Nous enregistrons donc le minimum strict.

\begin{theorem}[Théorème~\ref{thm:main}$'$]
\label{thm:main-prime}
Soient $M > 0$ et $G\colon[0, M] \to \R$ vérifiant
\[
p\,G(M) + (1 - p)\,G(0) \;=\; G(p\,M)
\qquad \text{pour tout } p \in [0, 1].
\]
Alors $G(v) = G(0) + \bigl(G(M) - G(0)\bigr)\,v/M$ sur~$[0, M]$.
\end{theorem}

La démonstration est mot pour mot celle du Théorème~\ref{thm:main}~: poser
$p = v/M$. Le Théorème~\ref{thm:main-prime} est en principe plus faible que
le Théorème~\ref{thm:main} --- l'hypothèse n'entraîne pas a
priori~\eqref{eq:star} dans sa totalité~; sous la conclusion (affinité),
les deux sont vérifiées. En pratique, un auteur qui
\emph{dérive}~\eqref{eq:star} d'un dispositif plus riche (par~ex.~d'une
identité d'égalité dans Jensen pour un calcul de risque bayésien, voir
Section~\ref{sec:recurrence}) dispose typiquement de~\eqref{eq:star} pour
tous les triplets $(u_1, u_2, p)$. Le Théorème~\ref{thm:main-prime}, plus
restreint, est la bonne référence pour un auteur qui souhaite mesurer la
part exacte de l'équation que la démonstration mobilise.

\subsection{Domaines convexes en dimensions supérieures}
\label{ssec:higher}

La substitution par corde s'étend mot pour mot à tout sous-ensemble
convexe~$C$ d'un espace vectoriel réel~$V$~: en appliquant le
Théorème~\ref{thm:main} le long de chaque corde, on montre que la
restriction de~$G$ à chaque segment de~$C$ est affine en le paramètre du
segment, et un argument inductif standard (cf.~\cite{AczelDhombres1989})
relève l'affinité-corde-par-corde en linéarité par combinaisons convexes
sur~$C$, donc en une forme affine $G(x) = a(x) + b$ pour une forme
linéaire~$a$ sur $\mathrm{vect}(C - C)$ et une constante $b \in \R$.
L'étape structurellement essentielle, en dimension~1, est le
Théorème~\ref{thm:main}.

\subsection{La version à coefficients rationnels $(J_{\Q})$ conserve,
elle, la pathologie}
\label{ssec:rational-pathology}

Le contraste $\Q$ vs $\R$ mérite d'être marqué explicitement.

\begin{proposition}[folklorique]
\label{prop:JQ-pathology}
Il existe $G\colon[0, 1] \to \R$ vérifiant~\eqref{eq:JQ} (et
donc~\eqref{eq:J2}) sur~$[0, 1]$, sans être affine.
\end{proposition}

\begin{proof}[Construction]
Choisissons une base de Hamel~$H$ de~$\R$ sur~$\Q$ contenant~$1$.
Définissons une application $\Q$-linéaire $\ell\colon\R \to \R$ en posant
$\ell(1) = 0$ et $\ell(h) = 1$ pour un certain élément de
base~$h \in H \setminus \{1\}$, et arbitrairement sur les autres éléments
de base. Par $\Q$-linéarité, $\ell(x + y) = \ell(x) + \ell(y)$ pour tous
$x, y \in \R$ (donc $\ell$ vérifie l'équation de Cauchy) et
$\ell(qx) = q\,\ell(x)$ pour tous $q \in \Q$ et $x \in \R$
($\Q$-homogénéité). Mais $\ell$ n'est \emph{pas} $\R$-linéaire~:
$\ell(1) = 0$ tandis que $\ell(h) = 1$, de sorte que les valeurs de~$\ell$
sur l'enveloppe $\Q$-linéaire de~$\{1\}$ (à savoir~$\Q$, sur laquelle
$\ell \equiv 0$) et sur l'enveloppe $\Q$-linéaire de~$\{h\}$ (où $\ell$
est non nulle) sont incompatibles avec toute forme $\R$-linéaire. Posons
$G := \ell|_{[0, 1]}$. Alors $G$ vérifie~\eqref{eq:JQ} sur~$[0, 1]$ ---
car les combinaisons convexes à coefficients rationnels de points
de~$[0, 1]$ restent dans~$[0, 1]$ et $\ell$ respecte les combinaisons
$\Q$-linéaires sur~$\R$ entier.

Témoin explicite de la non-affinité de~$G$~: choisissons
$h \in H \setminus \{1\}$ avec $h \in (0, 1)$ --- un tel élément de base
existe, l'extension de tout irrationnel de~$(0, 1)$ en une base de Hamel
étant immédiate. Alors $G(h) = \ell(h) = 1 \neq 0$, tandis que
$G(q) = \ell(q) = q\,\ell(1) = 0$ pour tout rationnel $q \in [0, 1]$.
Toute application affine $A\colon[0, 1] \to \R$ coïncidant avec~$G$ sur
$\Q \cap [0, 1]$ vérifierait $A(0) = 0$ et $A(q) = 0$ pour tout rationnel
$q \in [0, 1]$, ce qui force $A \equiv 0$ --- en contradiction avec
$G(h) = 1$. Donc $G$ n'est pas affine.
\end{proof}

La Proposition~\ref{prop:JQ-pathology} affine l'énoncé du
Théorème~\ref{thm:main} --- \emph{« aucune hypothèse de régularité n'est
requise »} --- en exhibant que le renforcement, du rationnel au continu
dans~\eqref{eq:star}, accomplit un travail effectif. Sans ce renforcement
(i.e., pour~\eqref{eq:JQ}), des hypothèses de régularité sont véritablement
nécessaires~; avec lui, elles sont vestigiales.

% =============================================================================
\section{Où le piège ressurgit}
\label{sec:recurrence}

La substitution par corde est --- comme nous l'avons noté --- folklorique.
Les hypothèses de régularité qu'elle rend vestigiales reparaissent
néanmoins dans les travaux appliqués qui dérivent~\eqref{eq:star} en
dehors de la communauté des équations fonctionnelles. Nous documentons
\emph{une} telle occurrence ci-dessous --- dans nos propres travaux sur le
seuil d'erreur atteignable des classifieurs fondés sur des partitions ---
et articulons la \emph{source structurelle} sous laquelle l'occurrence est
prévisible dans toute dérivation de calibration suffisamment riche. Nous
invitons les lecteurs qui rencontrent~\eqref{eq:star} dans leurs propres
travaux à compléter ce catalogue.

\subsection{Une occurrence : calibration par perte de substitution sur
l'axe de résolution}
\label{ssec:recurrence}

Dans un travail récent~\cite{El1,El2}, l'équation~\eqref{eq:star} surgit
avec~$G$ la fonction qui exprime le \emph{risque bayésien partitionnel}
d'un classifieur mesurable en fonction d'une fonctionnelle de score
concave agrégée sur les cellules de la partition~: $u_1, u_2$ sont les
valeurs du score par cellule et~$p$ est la masse de la cellule, qui
parcourt librement~$[0, 1]$ sur un espace de probabilité sous-jacent non
atomique~; ainsi~\eqref{eq:star} est vérifiée dans toute sa généralité.
Dans la formalisation Lean~4 de~\cite{El2}, le lemme correspondant a
d'abord été déclaré avec une hypothèse de bornitude dans sa signature, par
déférence à la littérature de Cauchy/Hamel~; le corps de la démonstration
a ensuite mis en évidence que cette hypothèse n'était pas mobilisée,
précisément par la substitution par corde du Théorème~\ref{thm:main}. Les
versions antérieures du texte principal invoquaient en conséquence la
remarque d'excuse \emph{« $G$ est bornée, donc le théorème de
Jensen-borné-implique-affine s'applique »} --- remarque que le
Théorème~\ref{thm:main} de la présente note retire.

\subsection{Source structurelle : pourquoi le piège est appelé à se
répéter}
\label{ssec:structural}

L'occurrence est prévisible. Dès qu'un argument de théorie de la
calibration aboutit à une identité de la forme
\[
\E[g(\xi)] \;=\; g\bigl(\E[\xi]\bigr)
\qquad (\xi \text{ variable aléatoire à valeurs dans } I,\
g\colon I \to \R),
\]
\emph{pour une classe de variables aléatoires~$\xi$ assez large pour que
l'espérance marginale~$\E[\xi]$ puisse être en tout point de~$I$ et que
le support de~$\xi$ puisse être tout sous-ensemble à deux points de~$I$
avec n'importe quelle paire de masses $(p, 1-p)$ avec $p \in [0, 1]$},
l'identité est~\eqref{eq:star} avec $g = G$ et~$\xi$ supportée
par~$\{u_1, u_2\}$. Une telle identité apparaît dès que l'inégalité de
Jensen est poussée à saturation. La littérature standard de la calibration
par perte de substitution --- Bartlett--Jordan--McAuliffe~\cite{BJM2006},
Tewari--Bartlett~\cite{TewariBartlett2007},
Steinwart~\cite{Steinwart2007},
Reid--Williamson~\cite{ReidWilliamson2010, ReidWilliamson2011} --- dérive
ses résultats de calibration par voie d'analyse convexe (biconjuguées,
dualité de Fenchel, hyperplans d'appui), qui maintient Jensen comme
\emph{inégalité} avec un défaut explicitement contrôlé et, par conséquent,
ne produit jamais~\eqref{eq:star} comme égalité~; la dérivation par axe
de résolution de~\cite{El2} pousse Jensen jusqu'à l'égalité et produit
ainsi~\eqref{eq:star} directement. La substitution par corde clôt
l'équation résultante en une ligne. D'autres styles de dérivation qui
saturent l'inégalité de Jensen sur une classe assez large de distributions
bipoints devraient produire~\eqref{eq:star} de manière analogue ---
moment où le Théorème~\ref{thm:main} retire instantanément la question des
hypothèses de régularité.

\subsection{Cadres adjacents}
\label{ssec:adjacent}

Deux cadres adjacents où des équations structurellement similaires
surgissent et où la substitution par corde constitue la fermeture
naturelle, sans pour autant que le \emph{piège} (l'invocation inutile de
la régularité de Cauchy/Hamel) soit, à notre connaissance, récurrent dans
la littérature publiée. Premièrement, les \textbf{théorèmes de
représentation de l'utilité espérée} dans la tradition de
von~Neumann--Morgenstern dérivent la linéarité-en-probabilité d'une
fonctionnelle d'utilité $U(p\,L_1 + (1-p)\,L_2) = p\,U(L_1) + (1-p)\,U(L_2)$
sur les loteries~; l'axiomatisation de Herstein--Milnor clôt cette
équation par l'\emph{axiome d'Archimède} plutôt que par Cauchy/Hamel, mais
la substitution par corde est l'alternative naturelle. Deuxièmement, la
\textbf{caractérisation axiomatique de l'entropie de Shannon} (axiome de
récursivité de Khinchin--Faddeev~; cf.~\cite{Aczel1966}
et~\cite{AczelDhombres1989}) traite des équations fonctionnelles plus
riches que~\eqref{eq:star}, mais plusieurs étapes intermédiaires se
ramènent à des identités de type Jensen à coefficient continu pour
lesquelles la substitution par corde fait partie de la boîte à outils
standard. Nous signalons ces cadres comme adjacents plutôt que comme
instances documentées du piège.

\subsection{Une invitation à compléter}
\label{ssec:invitation}

Nous invitons les lecteurs ayant rencontré~\eqref{eq:star} dans leurs
propres travaux et ayant invoqué une hypothèse de régularité vestigiale à
compléter ce catalogue par une référence d'une ligne renvoyant à l'étape
pertinente de leur article.

% =============================================================================
\section*{Déclaration d'intérêts}

L'auteur ne travaille pas pour, ne conseille pas, ne possède pas de parts
dans, et ne reçoit pas de financements d'une organisation qui pourrait
bénéficier de cet article, et n'a déclaré aucune affiliation autre que son
organisme de recherche.

% =============================================================================
\section*{Remerciements}

L'observation par substitution par corde a émergé dans le volet de
formalisation Lean~4 de~\cite{El2}, durant une phase du développement de
la démonstration dont la discipline est consignée dans trois
compétences-projet publiquement documentées au sein du dépôt-source
associé à~\cite{El2}. L'auteur remercie le processus de relecture
adversariale interne, décrit dans la même source, pour avoir mis en
évidence l'hypothèse de bornitude sur-engagée dont le retrait a motivé la
présente note.

% =============================================================================
% Bibliographie
% =============================================================================

\bibliographystyle{plain} % REPLI : remplacer par le style `crmath` quand
                          % la classe officielle Centre Mersenne est installée
\bibliography{refs}

\end{document}
